\chapter{Dos átomos às galáxias: as escalas do Universo}\label{ch:g9:atoms-to-galaxies}

Nove anos atrás este livro começou com os seus cinco sentidos e as
patas de uma joaninha. Ele se fecha com uma viagem que sentido nenhum
consegue fazer sozinho: quarenta e uma potências de dez, do coração
de um \omterm{def:g9:atoms-to-galaxies:atom}{átomo} até a borda do universo observável --- com você, bem
exatamente, no meio. Um último capítulo, um último segredo devido de
uma velha piada, e uma escada onde se apoiar para toda a física que
está por vir.

\section{Para baixo: matéria adentro}

\begin{definition}[Átomos]\label{def:g9:atoms-to-galaxies:atom}
As \omterm{def:g7:states-of-matter:molecule}{moléculas} do seu sétimo ano são elas mesmas construídas: toda
\omterm{def:g7:states-of-matter:molecule}{molécula} é um conjunto de \emph{átomos}\index{átomo} --- o alfabeto
da natureza, com cerca de cem tipos, cujas combinações soletram todas
as substâncias. Uma \omterm{def:g7:states-of-matter:molecule}{molécula} de água: dois átomos de hidrogênio e um
de oxigênio. Um átomo, por sua vez, tem arquitetura: um
\emph{núcleo}\index{núcleo} minúsculo, pesado e de carga positiva no
centro e, bem longe dele, uma nuvem de
\emph{elétrons}\index{elétron} --- portadores quase sem \omterm{def:g4:weight-and-mass:weight}{peso} de carga
negativa. Os tamanhos dos átomos ficam perto de $10^{-10}$ \unit{m};
os núcleos deles, perto de $10^{-15}$ --- cem mil vezes menores
ainda.
\end{definition}

\begin{proposition}[A matéria é quase toda vazio]\label{prop:g9:atoms-to-galaxies:empty}
Amplie um \omterm{def:g9:atoms-to-galaxies:atom}{átomo} até o tamanho de um grande estádio, e o \omterm{def:g9:atoms-to-galaxies:atom}{núcleo} dele é
uma ervilha no círculo central --- os \omterm{def:g9:atoms-to-galaxies:atom}{elétrons}, mosquitinhos nas
arquibancadas mais altas: tudo o que há no meio é vazio. A mesa de
mármore, a bigorna de ferro, a sua própria mão: esmagadoramente
vazias, com a solidez delas sendo uma recusa elétrica --- as nuvens
de \omterm{def:g9:atoms-to-galaxies:atom}{elétrons} dos \omterm{def:g9:atoms-to-galaxies:atom}{átomos} se repelindo até criar a ilusão de plenitude.
É uma das frases mais espantosas da física, e é a aritmética simples
de $10^{-10}$ contra $10^{-15}$.
\end{proposition}

\begin{example}[A velha piada, paga]\label{ex:g9:atoms-to-galaxies:electron}
Os capítulos de eletricidade deviam a você uma piadinha, e aqui está
ela. As cargas em marcha em todo fio deste livro são
\emph{\omterm{def:g9:atoms-to-galaxies:atom}{elétrons}} --- e eles são negativos, de modo que marcham do
\omterm{def:g3:first-electric-circuit:battery}{terminal} $-$ da \omterm{def:g3:first-electric-circuit:battery}{pilha} para o $+$: precisamente ao contrário do
sentido convencional fixado pelos físicos um século antes de alguém
poder perguntar aos marchantes. A convenção foi mantida --- toda lei
dos seus circuitos funciona perfeitamente com ela ---, mas a
ironiazinha continua de pé: em todo esquema que você desenhou, a
multidão de verdade caminha para o outro lado.
\end{example}

\begin{example}[E mais para baixo]\label{ex:g9:atoms-to-galaxies:deeper}
O \omterm{def:g9:atoms-to-galaxies:atom}{núcleo} tem partes próprias, e a história delas --- como se ligam,
por que alguns \omterm{def:g9:atoms-to-galaxies:atom}{núcleos} se dividem (o \omterm{def:g5:heat-and-insulation:heat}{calor} da usina) ou se fundem (o
do Sol) --- leva a física até $10^{-15}$ \omterm{def:g2:measuring-length:units}{metros} e além, para dentro
das menores estruturas que os seres humanos sondaram. O volume do
ensino médio abre o \omterm{def:g9:atoms-to-galaxies:atom}{núcleo} como se deve; os anos de universidade vão
mais longe ainda. Os degraus mais baixos da escada continuam sendo
construídos.
\end{example}

\section{Para cima: céu adentro}

\begin{example}[A subida]\label{ex:g9:atoms-to-galaxies:ascent}
Suba agora, potência de dez por potência de dez, passando pelos
marcos familiares de nove anos. Você, uns $10^{0}$ \unit{m}. O pátio
da escola, $10^{2}$. A Terra, $10^{7}$ de ponta a ponta --- a bolinha
de gude do modelo do campo esportivo. A distância da \omterm{def:g2:moon-in-the-sky:moon}{Lua},
$4 \times 10^{8}$; a do Sol, $1.5 \times 10^{11}$ --- oito
\omterm{def:g1:light-and-shadows:source}{minutos-luz}. A largura do \omterm{def:g4:solar-system:family}{Sistema Solar}, mais ou menos $10^{13}$; a
\omterm{def:g4:solar-system:star}{estrela} mais próxima, $4 \times 10^{16}$ --- quatro \omterm{def:g1:light-and-shadows:source}{anos-luz}. A \omterm{ex:g5:stars-night-sky:milkyway}{Via
Láctea}, $10^{21}$ \omterm{def:g2:measuring-length:units}{metros} de cidade; a galáxia vizinha,
$2 \times 10^{22}$ de distância; e o céu mais profundo já mapeado,
perto de $10^{26}$ \omterm{def:g2:measuring-length:units}{metros} --- com a \omterm{def:g1:light-and-shadows:source}{luz} das galáxias dele mais velha
que a Terra debaixo da sua cadeira.
\end{example}

\begin{omfigure}
\begin{tikzpicture}[scale=0.85]
  \draw[-Stealth, very thick] (0,0.6) -- (11.4,0.6);
  \foreach \x/\p in {0.4/-15, 1.7/-10, 3.0/-5, 4.3/0, 5.6/5,
      6.9/10, 8.2/15, 9.5/20, 10.8/25} {
    \draw[thick] (\x,0.45) -- (\x,0.75);
    \node[below, font=\footnotesize] at (\x,0.4) {$10^{\p}$};
  }
  % labels on two alternating rows, with thin leaders
  \foreach \x/\h/\name in {0.4/0.95/núcleo, 1.7/1.6/átomo,
      3.0/0.95/formiga, 4.3/1.6/você, 5.6/0.95/{um país},
      6.9/1.6/{distância do Sol}, 8.2/0.95/{estrelas próximas},
      9.5/1.6/{a galáxia}, 10.8/0.95/{céu profundo}} {
    \draw[omExo!80, thin] (\x,0.75) -- (\x,\h);
    \node[above, font=\footnotesize] at (\x,\h) {\name};
  }
\end{tikzpicture}

{\small A escada das escalas, em \omterm{def:g2:measuring-length:units}{metros}: quarenta e uma potências de
dez do \omterm{def:g9:atoms-to-galaxies:atom}{núcleo} ao céu profundo --- com quem lê de pé quase exatamente
no degrau do meio. (A formiga perdoa os nossos arredondamentos.)}
\end{omfigure}

\begin{method}[Pensar em ordens de grandeza]\label{met:g9:atoms-to-galaxies:oom}
A primeira ferramenta do físico diante de qualquer pergunta nova:
\begin{enumerate}
  \item expresse a grandeza em notação científica;
  \item guarde apenas a potência de dez --- a \emph{ordem de
        grandeza}\index{ordem de grandeza};
  \item compare escadas, não algarismos: uma \omterm{def:g4:solar-system:star}{estrela} a $10^{16}$
        \omterm{def:g2:measuring-length:units}{metros} contra um \omterm{def:g4:solar-system:planet}{planeta} a $10^{11}$ difere por cinco
        degraus --- cem mil vezes, digam o que disserem os
        algarismos da frente;
  \item só então, se for preciso, afie os algarismos.
\end{enumerate}
Nove anos de passos de julgamento --- ``essa resposta faz sentido?''
--- amadurecem neste hábito. O primeiro capítulo do volume do ensino
médio constrói sobre ele a cultura de medida dele.
\end{method}

\begin{example}[Aritmética de escada]\label{ex:g9:atoms-to-galaxies:arithmetic}
Quantos \omterm{def:g9:atoms-to-galaxies:atom}{átomos} cabem na largura de um lápis? Lápis, $10^{-2}$
\unit{m}; \omterm{def:g9:atoms-to-galaxies:atom}{átomo}, $10^{-10}$: oito degraus --- cem milhões de \omterm{def:g9:atoms-to-galaxies:atom}{átomos},
ombro a ombro. Quantas Terras até o Sol? $1.5 \times 10^{11}$ sobre
$1.3 \times 10^{7}$: cerca de $10^{4}$ --- dez mil bolinhas de gude
do modelo do campo esportivo, enfileiradas. A escada responde em
segundos o que as contagens cruas de quilômetros enterram em zeros.
\end{example}

\section{A vista do degrau do meio}

\begin{remark}[O que nove anos construíram]\label{rem:g9:atoms-to-galaxies:built}
Olhe escada abaixo e conte o que é seu. Em $10^{-10}$: as \omterm{def:g7:states-of-matter:molecule}{moléculas} e
os \omterm{def:g9:atoms-to-galaxies:atom}{átomos} --- e com eles os três estados, o \omterm{def:g5:heat-and-insulation:heat}{calor}, o revezamento do
\omterm{def:g3:sound-around-us:sound}{som}, os marchantes da corrente. Em torno de $10^{0}$: as \omterm{def:g2:pushes-and-pulls:force}{forças} e a
\omterm{def:g5:everyday-energy:energy}{energia}, os raios retos da \omterm{def:g1:light-and-shadows:source}{luz}, as leis dos circuitos --- a física do
tamanho humano de nove anos de experiências. De $10^{7}$ a $10^{22}$:
a Terra que gira, as \omterm{def:g3:sun-days-seasons:year}{estações} inclinadas, a \omterm{def:g2:moon-in-the-sky:moon}{Lua} que cai, uma lei só
de \omterm{def:g9:gravitation:gravitation}{gravitação} comandando igualmente o pomar e a galáxia. Nenhuma
outra matéria escolar entrega a você uma única escada do interior da
matéria até a borda do visível --- e cada degrau foi escalado com
instrumentos, honestidade e uma aritmética que você mesmo pode
conferir.
\end{remark}

\begin{remark}[O que espera lá em cima]\label{rem:g9:atoms-to-galaxies:ahead}
E o que ficou por terminar é o melhor de tudo. Que lei, exatamente,
liga a \omterm{def:g2:pushes-and-pulls:force}{força} ao movimento? O que a \omterm{def:g1:light-and-shadows:source}{luz} \emph{é}, para correr a
$3 \times 10^{8}$ \omterm{def:g2:measuring-length:units}{metros} por segundo e carregar cores dentro de si? O
que se divide dentro do \omterm{def:g9:atoms-to-galaxies:atom}{núcleo} --- e o que forja a \omterm{def:g1:light-and-shadows:source}{luz} do Sol? Por
que a contagem da \omterm{def:g5:everyday-energy:energy}{energia} útil sempre corre ladeira abaixo? Cada
pergunta foi honestamente levantada e honestamente adiada nestas
páginas; cada uma tem resposta, e as respostas estão entre as coisas
mais finas que a nossa espécie conhece. Elas começam no volume do
ensino médio, daqui a um ano. Leve a escada --- e o hábito de
conferir tudo.
\end{remark}

\section{Exercícios}

\begin{exercise}[$\star$]\label{exo:g9:atoms-to-galaxies:1}
Ordene por tamanho: \omterm{def:g7:states-of-matter:molecule}{molécula}, \omterm{def:g9:atoms-to-galaxies:atom}{átomo}, \omterm{def:g9:atoms-to-galaxies:atom}{núcleo}, formiga. Dê a potência
de dez de cada um em \omterm{def:g2:measuring-length:units}{metros}.
\end{exercise}

\begin{exercise}[$\star$]\label{exo:g9:atoms-to-galaxies:2}
Descreva a arquitetura do \omterm{def:g9:atoms-to-galaxies:atom}{átomo} --- o que fica no centro, o que o
cerca e os dois tamanhos que fazem a \omterm{prop:g5:mirrors-reflection:image}{imagem} do estádio.
\end{exercise}

\begin{exercise}[$\star$]\label{exo:g9:atoms-to-galaxies:3}
Pague a velha piada adiante: quem marcha de verdade num fio, para que
lado --- e por que todas as leis dos circuitos sobrevivem à
revelação?
\end{exercise}

\begin{exercise}[$\star$]\label{exo:g9:atoms-to-galaxies:4}
O que é uma \omterm{met:g9:atoms-to-galaxies:oom}{ordem de grandeza}? Dê a \omterm{met:g9:atoms-to-galaxies:oom}{ordem de grandeza} de: a sua
altura; o comprimento da sala de aula; o diâmetro da Terra.
\end{exercise}

\begin{exercise}[$\star$]\label{exo:g9:atoms-to-galaxies:5}
Por quantos degraus da escada estes diferem: o \omterm{def:g9:atoms-to-galaxies:atom}{átomo} e o \omterm{def:g9:atoms-to-galaxies:atom}{núcleo}? você
e a Terra? a distância do Sol e a da \omterm{def:g4:solar-system:star}{estrela} mais próxima?
\end{exercise}

\begin{exercise}[$\star$]\label{exo:g9:atoms-to-galaxies:6}
``A mesa é quase toda espaço vazio.'' Defenda a frase com os dois
números atômicos --- e explique o que faz a mesa parecer sólida mesmo
assim.
\end{exercise}

\begin{exercise}[$\star$]\label{exo:g9:atoms-to-galaxies:7}
Aritmética de escada: mais ou menos quantos \omterm{def:g9:atoms-to-galaxies:atom}{átomos}, enfileirados,
cobrem um milímetro?
\end{exercise}

\begin{exercise}[$\star$]\label{exo:g9:atoms-to-galaxies:8}
Localize na escada, com uma potência de dez para cada: um \omterm{def:g8:speed-of-light:lightyear}{ano-luz}; a
\omterm{ex:g5:stars-night-sky:milkyway}{Via Láctea}; o céu mais profundo já mapeado.
\end{exercise}

\begin{exercise}[$\star\star$]\label{exo:g9:atoms-to-galaxies:9}
O modelo do estádio: se o \omterm{def:g9:atoms-to-galaxies:atom}{núcleo} é uma ervilha de \qty{1}{cm}, qual é
a largura do \omterm{def:g9:atoms-to-galaxies:atom}{estádio-átomo}? (Razão de $10^{5}$ --- converta para
\omterm{def:g2:measuring-length:units}{metros} e julgue contra um estádio de verdade.)
\end{exercise}

\begin{exercise}[$\star\star$]\label{exo:g9:atoms-to-galaxies:10}
Uma gota de água guarda cerca de $10^{21}$ \omterm{def:g7:states-of-matter:molecule}{moléculas}. Usando a
comparação do mar de gotas do capítulo do seu sétimo ano --- ou a sua
própria estimativa das gotas de todos os oceanos (uns $10^{21}$
\omterm{def:g6:mass-and-volume:volume}{litros}, com $10^{5}$ gotas cada um) --- pese a velha afirmação de que
uma gota guarda mais \omterm{def:g7:states-of-matter:molecule}{moléculas} do que os mares guardam gotas.
\end{exercise}

\begin{exercise}[$\star\star$]\label{exo:g9:atoms-to-galaxies:11}
O seu corpo fica perto do meio da escada: mais ou menos quantos
degraus para baixo até o \omterm{def:g9:atoms-to-galaxies:atom}{átomo}, e quantos para cima até o céu mais
profundo? O que essa simetria diz sobre o alcance de nove anos de
escola?
\end{exercise}

\begin{exercise}[$\star\star\star$]\label{exo:g9:atoms-to-galaxies:12}
O último exercício do livro. Escolha três capítulos quaisquer dos
seus nove anos --- um dos anos 1 a 3, um dos anos 4 a 6, um dos anos
7 a 9 --- e escreva, para cada um, a única frase que você diria a si
mesmo mais novo, começando aquele capítulo: o que ele de fato
ensinou, visto do alto da escada. (Não há página de solução para uma
carta a si mesmo --- mas escreva; o curso que vem pela frente vai
ficar feliz que você tenha escrito.)
\end{exercise}

\section{Problema: A viagem em potências de dez}

\begin{problem}[{Problema de fim de semana --- a turma filma ``A
viagem em potências de dez''; um zoom, quarenta e um degraus; a
reunião de roteiro}]\label{pb:g9:atoms-to-galaxies:1}
O filme da turma: um zoom contínuo de uma toalha de piquenique para
fora, até o céu profundo, e depois para dentro, até o \omterm{def:g9:atoms-to-galaxies:atom}{núcleo} --- uma
potência de dez por segundo de filme. Você escreve o roteiro
científico.

\medskip\noindent\textbf{Parte I --- Para fora.}
O zoom começa enquadrando um \omterm{def:g2:measuring-length:units}{metro} de toalha: $10^{0}$.
\begin{enumerate}
  \item Em que segundos o quadro contém pela primeira vez: o pátio
        inteiro da escola ($10^{2}$); a Terra inteira ($10^{7}$); a
        \omterm{def:g6:sun-earth-moon:axisorbit}{órbita} da \omterm{def:g2:moon-in-the-sky:moon}{Lua} ($10^{9}$)?
  \item Que segundo enquadra a distância do Sol --- e o que o
        narrador diz sobre a idade da \omterm{def:g1:light-and-shadows:source}{luz} do Sol, tirado do capítulo
        da \omterm{def:g1:light-and-shadows:source}{luz}?
  \item Entre que segundos o quadro atravessa o grande vazio em que o
        \omterm{def:g4:solar-system:family}{Sistema Solar} já acabou mas a \omterm{def:g4:solar-system:star}{estrela} mais próxima ainda não
        entrou? O que a tela deveria honestamente mostrar ali?
  \item O zoom termina no segundo $26$: o que preenche o quadro, e
        que idade tem a \omterm{def:g1:light-and-shadows:source}{luz} mais antiga dele em comparação com a
        Terra?
\end{enumerate}

\medskip\noindent\textbf{Parte II --- Para dentro.}
O zoom se inverte e mergulha numa folha caída sobre a toalha.
\begin{enumerate}[resume]
  \item Em que segundos o filme passa por: uma formiga ($10^{-2}$,
        com generosidade); uma célula ($10^{-5}$); uma \omterm{def:g7:states-of-matter:molecule}{molécula}
        ($10^{-9}$); um \omterm{def:g9:atoms-to-galaxies:atom}{átomo} ($10^{-10}$)?
  \item Entre a nuvem de \omterm{def:g9:atoms-to-galaxies:atom}{elétrons} do \omterm{def:g9:atoms-to-galaxies:atom}{átomo} e o \omterm{def:g9:atoms-to-galaxies:atom}{núcleo} dele, o filme
        precisa dar zoom por mais cinco segundos através de --- o
        quê? O que o narrador diz sobre o estádio?
  \item O quadro final do filme, em $10^{-15}$: o assunto dele --- e
        a honesta frase de encerramento do narrador sobre os degraus
        de baixo.
  \item Some a \omterm{def:g2:measuring-time:duration}{duração} do filme: quantos segundos para fora, quantos
        para dentro, a partir do \omterm{def:g2:measuring-length:units}{metro} da toalha?
\end{enumerate}

\medskip\noindent\textbf{Parte III --- A reunião de roteiro.}
\begin{enumerate}[resume]
  \item Um professor objeta que o filme gasta um segundo no passo da
        casa até a rua e um segundo no passo da galáxia até o
        aglomerado de galáxias --- ``passos absurdamente desiguais!''.
        Defenda os segundos iguais com a lógica da escada.
  \item O narrador quer uma frase recorrente para as duas direções
        --- algo verdadeiro em cada degrau a respeito do vazio.
        Rascunhe-a (o \omterm{def:g4:solar-system:family}{Sistema Solar} e o \omterm{def:g9:atoms-to-galaxies:atom}{átomo} precisam caber os dois
        embaixo dela).
  \item Os créditos precisam dizer o nome dos dois instrumentos que
        conquistaram as duas direções. Diga o nome deles e os
        capítulos deste livro que apresentaram os princípios deles.
  \item As últimas palavras do filme pertencem ao próprio curso:
        escreva o epílogo de duas frases --- nove anos, uma escada e
        o que começa no ano que vem.
\end{enumerate}
\end{problem}
